% !PRINTED 376
\weberpage{376}
%<<CONT
constanten Factor in $\xi_{41}$ verfügt wird, $f$ auch in der Form darstellen:

\weq{7}{f = 4x_2\xi_3\,x_4\xi_{41} - v^2,}

worin $v$ eine quadratische Form ist. Hieraus und aus der ersten Darstellung in (5) ergiebt sich aber die Identität

\[
4x_2\xi_3(x_3\xi_2 - x_4\xi_{41}) = (u_1 - v)(u_1 + v),
\]

und daraus schliessen wir, dass einer der beiden Factoren rechts durch $x_2\xi_3$ theilbar sein muss [\S.\,96, (3)]. Nehmen wir an, es sei dies $u_1 - v$, was durch Verfügung über das Vorzeichen von $v$ erreicht werden kann, so folgt, dass ein von Null verschiedener constanter Factor $\lambda$ existiren muss, so dass
% [sic] The source prints \lambda here; the following display and later prose use \lambda_1.

\[
\begin{aligned}
u_1 - v &= 2\lambda_1 x_2\xi_3,\\
u_1 + v &= \frac{2(x_3\xi_2 - x_4\xi_{41})}{\lambda_1},
\end{aligned}
\]

woraus

\[
u_1 = \lambda_1 x_2\xi_3 + \frac{x_3\xi_2 - x_4\xi_{41}}{\lambda_1}
\]

folgt. Ersetzen wir hierin $u_1$ durch seinen Ausdruck aus (6), so ergiebt sich

\[
x_4\xi_{41} = x_3\xi_2 - \lambda_1(-x_1\xi_1 + x_2\xi_2 + x_3\xi_3) + \lambda_1^2 x_2\xi_3.
\]

Solcher Gleichungen erhalten wir zunächst drei, wenn wir die Indices 1, 2, 3 cyklisch vertauschen und an Stelle der unbekannten Constanten $\lambda_1$ drei Constanten $\lambda_1, \lambda_2, \lambda_3$ setzen:

\weq{8}{\begin{aligned}
x_4\xi_{41} &= x_3\xi_2 - \lambda_1(-x_1\xi_1 + x_2\xi_2 + x_3\xi_3) + \lambda_1^2 x_2\xi_3\\
x_4\xi_{42} &= x_1\xi_3 - \lambda_2(x_1\xi_1 - x_2\xi_2 + x_3\xi_3) + \lambda_2^2 x_3\xi_1\\
x_4\xi_{43} &= x_2\xi_1 - \lambda_3(x_1\xi_1 + x_2\xi_2 - x_3\xi_3) + \lambda_3^2 x_1\xi_2.
\end{aligned}}

Um die Constanten $\lambda_1, \lambda_2, \lambda_3$ zu bestimmen, dividiren wir die beiden letzten dieser Gleichungen mit $\lambda_2$ und $\lambda_3$ und addiren. Dadurch folgt die identische Gleichung

\weq{9}{\begin{aligned}
x_4\left(\frac{\xi_{42}}{\lambda_2} + \frac{\xi_{43}}{\lambda_3}\right) &=\\
x_1\left(-2\xi_1 + \lambda_3\xi_2 + \frac{\xi_3}{\lambda_2}\right) &+ \xi_1\left(\lambda_2 x_3 + \frac{x_2}{\lambda_3}\right).
\end{aligned}}

Da nun $x_4, x_1, \xi_1$ azygetisch sind, und folglich nicht durch einen Punkt gehen können, so muss die Linie

\[
\lambda_2 x_3 + \frac{x_2}{\lambda_3} = 0
\]
%CONT>>
